\documentclass[11pt,a4paper]{article}
\usepackage[T1]{fontenc}
\usepackage{lmodern}
\usepackage[margin=28mm]{geometry}
\usepackage{amsmath,amssymb,amsthm,booktabs}
\usepackage[colorlinks=true,linkcolor=blue,citecolor=blue,urlcolor=blue]{hyperref}
\hypersetup{pdftitle={A four-dimensional periodic Costas array of order six},pdfauthor={Kaiwen Jin}}
\newtheorem{theorem}{Theorem}
\newcommand{\Z}{\mathbb Z}
\newcommand{\res}{\operatorname{res}}
\title{A four-dimensional periodic Costas array of order six}
\author{Kaiwen Jin\\\small\texttt{yui.ui.siki@gmail.com}}
\date{15 October 2026}
\begin{document}
\maketitle
\begin{abstract}
We give a bijection from $[2]\times[3]$ onto itself whose graph, extended
periodically in four coordinates, has distinct nonzero ordered ordinary
difference vectors in every $2\times3\times2\times3$ window. Its order is
$6$, whereas the proposed minimum-order restriction of Rubio and Torres
would require $4$. A short proof keeps the signs of the two binary-coordinate
differences and reduces the other two differences modulo three.
\end{abstract}

\section{The periodicity question}

Write $[r]=\{1,\ldots,r\}$. Rubio and Torres~\cite{RT} define a
multidimensional permutation array through a bijection
\[
 \varphi:\prod_{i=1}^k[n_i]\longrightarrow\prod_{i=k+1}^m[n_i],
 \qquad n=\prod_{i=1}^k n_i=\prod_{i=k+1}^m n_i,
\]
where $m\geq2$ and all side lengths $n_i\geq2$. If $S$ is its graph and
$L=n_1\Z\times\cdots\times n_m\Z$, the array is periodic Costas when,
for every $t\in\Z^m$, the nonzero ordered differences between points of
\[
 (S+L)\cap\left(t+\prod_{i=1}^m[n_i]\right)
\]
are distinct. These differences are ordinary integer vectors. In
Conjecture~1 of~\cite{RT}, Rubio and Torres propose that, when
$1\leq m-k\leq k$, periodicity forces $n=2^k$, or equivalently
$n_1=\cdots=n_k=2$.

The six-point array below contradicts this restriction in dimension four.
An earlier public Prove2Me record by shivm, dated 4~October~2026, labels
the same conjecture disproved~\cite{prior}. We have not inspected the proof
or witness underlying that record. We acknowledge the earlier announcement
and make no novelty or priority claim. The argument here is self-contained.

\clearpage
\section{An array of order six}

Let $\varphi:[2]\times[3]\to[2]\times[3]$ be the map in Table~\ref{tab:array},
and let $P_1,\ldots,P_6$ denote the corresponding graph points.

\begin{table}[ht]
\centering
\begin{tabular}{ccc}
\toprule
$x$ & $\varphi(x)$ & $(x,\varphi(x))$\\
\midrule
$(1,1)$ & $(1,1)$ & $P_1=(1,1,1,1)$\\
$(1,2)$ & $(1,2)$ & $P_2=(1,2,1,2)$\\
$(1,3)$ & $(2,1)$ & $P_3=(1,3,2,1)$\\
$(2,1)$ & $(1,3)$ & $P_4=(2,1,1,3)$\\
$(2,2)$ & $(2,3)$ & $P_5=(2,2,2,3)$\\
$(2,3)$ & $(2,2)$ & $P_6=(2,3,2,2)$\\
\bottomrule
\end{tabular}
\caption{The bijection defining the array.}
\label{tab:array}
\end{table}

\begin{theorem}\label{thm:main}
Set $S=\{P_1,\ldots,P_6\}$ and $L=2\Z\times3\Z\times2\Z\times3\Z$.
For every $t\in\Z^4$, all nonzero ordered ordinary differences between
points of
\[
 S_t=(S+L)\cap\bigl(t+[2]\times[3]\times[2]\times[3]\bigr)
\]
are distinct. Thus $S$ is a periodic four-dimensional Costas array of
order $6$ with $k=2$, and $6>2^2$.
\end{theorem}

\begin{proof}
The input and output columns of Table~\ref{tab:array} each enumerate
$[2]\times[3]$, so $\varphi$ is a bijection. To distinguish its ordered
displacements, retain the two binary-coordinate signs. For
$P=(a,b,c,d)$ and $Q=(a',b',c',d')$ in $S$, define
\begin{equation}\label{eq:H}
 H(P,Q)=\bigl(a'-a,\res_3(b'-b),c'-c,\res_3(d'-d)\bigr),
\end{equation}
where $\res_N(z)$ is the unique remainder in $\{0,\ldots,N-1\}$ of
the integer $z$ modulo $N$.

Table~\ref{tab:diff} gives the residue pairs in coordinates two and four,
grouped by the ordinary differences in coordinates one and three.
The ordered pair $(i,j)$ in the middle column means $(P_i,P_j)$.

\begin{table}[ht]
\centering
\small
\setlength{\tabcolsep}{5pt}
\begin{tabular}{ccc}
\toprule
$(a'-a,c'-c)$ & Ordered pairs & $(\res_3(b'-b),\res_3(d'-d))$\\
\midrule
$(0,0)$ & $(1,2),(2,1),(5,6),(6,5)$ & $(1,1),(2,2),(1,2),(2,1)$\\
$(0,1)$ & $(1,3),(2,3),(4,5),(4,6)$ & $(2,0),(1,2),(1,0),(2,2)$\\
$(1,0)$ & $(1,4),(2,4),(3,5),(3,6)$ & $(0,2),(2,1),(2,2),(0,1)$\\
$(1,1)$ & $(1,5),(1,6),(2,5),(2,6)$ & $(1,2),(2,1),(0,1),(1,0)$\\
$(1,-1)$ & $(3,4)$ & $(1,2)$\\
\bottomrule
\end{tabular}
\caption{Distinct residues within each fixed pair of binary differences.}
\label{tab:diff}
\end{table}

Reversing the pairs in the last four rows gives the remaining binary
differences
\[
 (0,-1),\quad(-1,0),\quad(-1,-1),\quad(-1,1)
\]
and negates the two residues
modulo three. Each row has distinct residue pairs, and different rows
have different binary differences. The displayed pairs and these reversals
are distinct and number $4+2(4+4+4+1)=30=6(6-1)$, so they exhaust
all ordered pairs of distinct points of $S$. Consequently $H$ is injective
on those pairs.

Now fix an arbitrary $t\in\Z^4$ and write $(n_1,n_2,n_3,n_4)=(2,3,2,3)$.
The representative of $u\in[n_i]$ in the interval $t_i+[n_i]$ is
\[
 F_{i,t}(u)=t_i+1+\res_{n_i}(u-t_i-1).
\]
It is congruent to $u$ modulo $n_i$ and lies in the stated interval.
Two congruent integers in that interval differ by a multiple of $n_i$
of absolute value at most $n_i-1$, so they are equal. It follows that
\[
 P\longmapsto P_t=\bigl(F_{1,t}(p_1),\ldots,F_{4,t}(p_4)\bigr)
\]
is a bijection from $S$ onto $S_t$.

For $i=1,3$, the values $F_{i,t}(1)$ and $F_{i,t}(2)$ are the two
consecutive integers in $t_i+[2]$. Hence there is a fixed
$\varepsilon_i\in\{-1,1\}$, depending only on $t_i$, such that
\[
 F_{i,t}(u')-F_{i,t}(u)=\varepsilon_i(u'-u),
 \qquad u,u'\in[2].
\]
For $i=2,4$, subtraction of congruences gives
\[
 \res_3\bigl(F_{i,t}(u')-F_{i,t}(u)\bigr)=\res_3(u'-u).
\]

Suppose $Q_t-P_t=V_t-U_t$ as ordinary integer vectors, where
$P,Q,U,V\in S$, $P\ne Q$ and $U\ne V$. The same signs
$\varepsilon_1,\varepsilon_3$ occur for both pairs. Cancelling them in
coordinates one and three, and reducing coordinates two and four modulo
three, gives $H(P,Q)=H(U,V)$. Its injectivity implies $P=U$ and $Q=V$,
and therefore $P_t=U_t$ and $Q_t=V_t$. Since $t$ was arbitrary, the required
distinct-difference property holds in every translated window.
\end{proof}

All four side lengths are at least two, $m-k=k=2$, and both products of
side lengths are $6$. Thus Theorem~\ref{thm:main} satisfies every premise
of the proposed restriction while its conclusion fails.

\paragraph{AI use.}
OpenAI Codex assisted in finding the six-point example, developing the proof,
drafting this note, and performing auxiliary exact-integer checks.
Separate fresh-context Codex passes assessed readability and challenged the
manuscript. Independent mathematical review is requested.

\begin{thebibliography}{9}
\raggedright
\bibitem{RT}
I. Rubio and J. Torres,
\emph{Multidimensional Costas Arrays and Their Periodicity},
\href{https://arxiv.org/abs/2208.02378v1}{arXiv:2208.02378v1} (2022).
Definitions 2--3 (p.~6), Definition 4 (p.~8), and Conjecture 1 (p.~16).

\bibitem{prior}
shivm,
\emph{Rubio--Torres Conjecture 1: periodic Costas arrays have order $2^k$},
Prove2Me, record dated 4 October 2026, marked ``Disproved''.
\url{https://prove2.me/theorems/f2732909-8aec-45d0-afe0-8943f5d2c80e}.
Accessed 15 October 2026; cited for the announcement.
\end{thebibliography}
\end{document}
